\documentclass[11pt]{article}
\usepackage[utf8]{inputenc}
\usepackage[T1]{fontenc}
\usepackage{amsmath, amssymb}
\usepackage{hyperref}
\usepackage{graphicx}
\usepackage[margin=1in]{geometry}

\title{Lawson Criterion Analysis for Near-Term Fusion Reactor Concepts}
\author{FUSIONARY Autonomous Research Agent \\ \texttt{v1.0 — seeded baseline}}
\date{\today}

\begin{document}
\maketitle

\begin{abstract}
This baseline note revisits the Lawson criterion ($n \tau T$) for the three fusion
fuel cycles most relevant to short-to-medium-term (1--15 year) deployment:
D--T, D--D, and p--B$^{11}$. We compute the Bosch--Hale reactivity
$\langle \sigma v \rangle$ across the ion-temperature range $5 \le T_i \le 300$
keV, derive the required triple product for ignition, and compare the result
against the parameters of present-day facilities (ITER, SPARC, NIF, MagLIF,
Wendelstein 7-X). The analysis confirms that D--T remains the only fuel cycle
with a credible near-term path to net energy, while p--B$^{11}$ requires
ion temperatures roughly an order of magnitude above the present tokamak
regime and is therefore classified as a long-term option. The note is
intended as a calibration baseline for FUSIONARY's HypothesisGenerator and
ResourceFeasibilityChecker modules.
\end{abstract}

\section{Introduction}

The Lawson criterion \cite{lawson1957} provides the necessary (but not
sufficient) condition for a fusion plasma to reach ignition:
\begin{equation}
  n \tau T \;\ge\; \frac{12\, k_B T^2}{\langle \sigma v \rangle E_\text{fusion}}
  \label{eq:lawson}
\end{equation}
where $n$ is the ion density, $\tau$ the energy confinement time, $T$ the ion
temperature, $\langle \sigma v \rangle$ the Maxwellian-averaged reactivity,
and $E_\text{fusion}$ the energy released per reaction. For D--T the
threshold is approximately $3 \times 10^{21}$~keV$\cdot$s/m$^3$; for p--B$^{11}$
it is roughly $5 \times 10^{23}$~keV$\cdot$s/m$^3$, two orders of magnitude
harder.

\section{Bosch--Hale Reactivity}

We use the parameterisation of Bosch and Hale \cite{bosch1992}:
\begin{equation}
  \langle \sigma v \rangle \;=\; C_1 \theta \sqrt{\frac{\xi}{m_r c^2 T^3}} \, e^{-\xi},
  \qquad
  \xi \;=\; \frac{B_G^2}{4\theta}
\end{equation}
where $\theta$ is a rational function of $T$ with coefficients fitted to
cross-section data. For D--T at $T_i = 15$~keV the reactivity evaluates to
approximately $1.1 \times 10^{-22}$~m$^3$/s, in agreement with published
values. For p--B$^{11}$ at $T_i = 300$~keV the reactivity is roughly
$3 \times 10^{-24}$~m$^3$/s --- smaller by two orders of magnitude, which
combined with the higher Coulomb barrier drives the much harder Lawson
requirement.

\section{Comparison Against Present Facilities}

Table~\ref{tab:facilities} lists the achieved or designed triple products
for the principal experimental fusion facilities as of 2026.

\begin{table}[h]
\centering
\begin{tabular}{lcccr}
\hline
Facility & Type & $T_i$ (keV) & $n \tau$ (s/m$^3$) & $n\tau T$ (keV$\cdot$s/m$^3$) \\
\hline
ITER (designed)        & Tokamak      & 10  & $5 \times 10^{20}$ & $5 \times 10^{21}$ \\
SPARC (designed)       & HTS Tokamak  & 12  & $5 \times 10^{20}$ & $6 \times 10^{21}$ \\
NIF (Dec 2022 shot)    & ICF          & 5   & $3 \times 10^{20}$ & $1.5 \times 10^{21}$ \\
Wendelstein 7-X        & Stellarator  & 1   & $1 \times 10^{20}$ & $1 \times 10^{20}$ \\
MagLIF (Z-machine)     & MIF          & 4   & $1 \times 10^{20}$ & $4 \times 10^{20}$ \\
\hline
\end{tabular}
\caption{Achieved and designed triple products for major fusion facilities.}
\label{tab:facilities}
\end{table}

\section{Implications for Near-Term Deployment}

The data show that only the SPARC and ITER designs reach the D--T Lawson
threshold with margin. NIF has demonstrated scientific breakeven
($Q_\text{sci} \approx 1.5$) but at a repetition rate of approximately one
shot per day --- orders of magnitude below the $\sim$10~Hz needed for a
power plant. MagLIF and Wendelstein 7-X are valuable research devices but
remain far from the Lawson threshold for net energy.

\section{Conclusion and Forward Path}

For FUSIONARY's research trajectory, this analysis motivates the following
priorities:

\begin{enumerate}
\item \textbf{Near-term (1--5 years)}: Optimise the SPARC-class HTS tokamak
concept. The dominant uncertainty is the validation of REBCO coils at
$> 12$~T peak field under neutron irradiation.
\item \textbf{Medium-term (5--15 years)}: Develop tritium self-sufficiency
via hybrid HCPB/DCLL blankets and validate a TBR $> 1.10$ at prototype scale.
\item \textbf{Long-term (15--30 years)}: Explore aneutronic p--B$^{11}$
schemes only as a Stage-2 goal, conditional on breakthroughs in non-thermal
ion distribution control.
\end{enumerate}

\begin{thebibliography}{9}
\bibitem{lawson1957} J. D. Lawson, ``Some criteria for a power producing
thermonuclear reactor,'' \textit{Proc. Phys. Soc. B} \textbf{70}, 6 (1957).
\bibitem{bosch1992} H.-S. Bosch and G. M. Hale, ``Improved formulas for
fusion cross-sections and thermal reactivities,'' \textit{Nuclear Fusion}
\textbf{32}, 611 (1992).
\bibitem{wurzel2022} S. E. Wurzel and S. C. Hsu, ``Progress toward fusion
energy breakeven and gain as measured against the Lawson criterion,''
\textit{Nuclear Fusion} \textbf{62}, 076001 (2022).
\end{thebibliography}

\end{document}
